\chapter{Espaços Euclidianos}\label{ch:b1:euclid}

Acrescentar um \omterm{def:b1:euclid:def}{produto interno} a um \omterm{def:b1:vspaces:def}{espaço vetorial} real compra as
noções geométricas --- comprimentos, ângulos, ortogonalidade,
distâncias --- e um teorema que domina o capítulo: todo \omterm{def:b1:vspaces:subspace}{subespaço}
admite uma \omterm{thm:b1:euclid:projection}{projeção ortogonal}, calculável por Gram--Schmidt, que
realiza a menor distância. As \omterm{def:b1:euclid:isometry}{isometrias} do plano fecham o capítulo
e a geometria do ano.

Ao longo do capítulo, $E$ é um \omterm{def:b1:vspaces:def}{espaço vetorial} \emph{real}.

\section{Produtos internos}

\begin{definition}\label{def:b1:euclid:def}
Um \emph{produto interno}\index{produto interno} em $E$ é uma
\omterm{def:b1:logic:map}{aplicação} $\langle\cdot,\cdot\rangle \colon E \times E \to \R$ bilinear, simétrica e definida positiva ($\langle x, x\rangle >
0$ para
$x \neq 0$). Um espaço de \omterm{def:b1:findim:def}{dimensão finita} assim munido é um \emph{espaço
euclidiano}\index{espaço euclidiano}. A \emph{norma} de $x$ é $\norm{x} = \sqrt{\langle x, x\rangle}$,
e $d(x, y) =
\norm{x - y}$.
\end{definition}

\begin{example}\label{ex:b1:euclid:examples}
Em $\R^n$: o produto canônico $\langle x, y\rangle = \sum x_i
y_i$. Em $C(\intcc{a}{b})$: $\langle f, g \rangle = \int_a^b fg$ (a definição
positiva é o \cref{thm:b1:integration:props}~(4)). Em $\R_n[X]$: $\langle P, Q\rangle = \int_0^1 PQ$, ou $\sum_{i}
P(x_i)Q(x_i)$ sobre $n+1$
pontos distintos.
\end{example}

\begin{example}[O ângulo entre dois polinômios]
\label{ex:b1:euclid:polyangle}
Uma vez escolhido um \omterm{def:b1:euclid:def}{produto interno}, \emph{quaisquer} dois vetores
não nulos têm um ângulo, via $\cos\theta = \frac{\langle x,
y\rangle}{\norm x\,\norm y}$ (um cosseno legítimo, por
Cauchy--Schwarz). Para $X$ e $X^2$ em $\int_0^1$:
\[
\langle X, X^2\rangle = \frac14, \qquad
\norm X = \frac1{\sqrt3}, \qquad \norm{X^2} = \frac1{\sqrt5},
\qquad
\cos\theta = \frac{1/4}{1/\sqrt{15}} = \frac{\sqrt{15}}{4}
\approx 0.968 :
\]
um ângulo de cerca de $14.5$ graus --- em $\intcc{0}{1}$, os gráficos de $x$
e $x^2$ são ``quase paralelos'' no sentido da média quadrática, e é
por isso que remover essa direção comum (Gram--Schmidt, abaixo)
deixa apenas a pequena correção $X^2 - X + \frac16$.
\end{example}

\begin{theorem}[Cauchy--Schwarz; propriedades da norma]\label{thm:b1:euclid:cs}
Para todos $x, y \in E$:\index{desigualdade de Cauchy--Schwarz}
\[
\abs{\langle x, y\rangle} \leq \norm x\, \norm y ,
\]
com igualdade se e somente se $x, y$ são proporcionais. Por
consequência, $\norm\cdot$ satisfaz a desigualdade triangular $\norm{x + y} \leq \norm x + \norm
y$ (e
$\norm{\lambda x} = \abs\lambda \norm x$, $\norm x = 0 \iff
x = 0$). Além disso:
\[
\norm{x+y}^2 = \norm x^2 + 2\langle x, y\rangle + \norm y^2,
\qquad
\langle x, y \rangle = \tfrac14\bigl(\norm{x+y}^2 -
\norm{x-y}^2\bigr) .
\]
\end{theorem}

\begin{proof}
Se $y = 0$, tudo é trivial. Caso contrário, a quadrática $t
\mapsto \norm{x + ty}^2 = \norm x^2 + 2t\langle x, y\rangle + t^2
\norm y^2$ é
$\geq 0$ para todo $t$: o seu discriminante é $\leq 0$, o que é
Cauchy--Schwarz; a igualdade significa uma raiz dupla $t_0$, isto é,
$x + t_0 y = 0$ (pela definição positiva): proporcionalidade. Desigualdade
triangular: desenvolva,
\[
\norm{x + y}^2 = \norm x^2 + 2\langle x, y\rangle + \norm y^2
\leq \norm x^2 + 2\norm x\,\norm y + \norm y^2
= \bigl(\norm x + \norm y\bigr)^2 ,
\]
sendo o passo do meio Cauchy--Schwarz; a igualdade força $\langle
x, y\rangle = \norm x\norm y$, o
caso de igualdade \emph{positivo}, isto é, proporcionalidade com
razão não negativa --- geometricamente, o triângulo só degenera
quando os dois vetores apontam no mesmo sentido. As duas últimas
identidades são desenvolvimentos diretos (a segunda, a
\emph{identidade de polarização}, recupera o produto a partir da
norma).
\end{proof}

\section{Ortogonalidade}

\begin{definition}\label{def:b1:euclid:orthogonal}
$x \perp y$ quando $\langle x, y \rangle = 0$. Uma família é \emph{ortogonal} quando os seus
vetores são dois a dois ortogonais, e
\emph{ortonormal}\index{família ortonormal} quando, além disso, cada
um tem norma $1$. O \emph{complemento ortogonal} de um \omterm{def:b1:vspaces:subspace}{subespaço} $F$
é
\[
F^{\perp} = \{x \in E : \forall y \in F,\ \langle x, y\rangle =
0\},
\]
um \omterm{def:b1:vspaces:subspace}{subespaço} de $E$.\index{complemento ortogonal}
\end{definition}

\begin{proposition}\label{prop:b1:euclid:orthfree}
(Pitágoras\index{teorema de Pitágoras}) Se $x \perp y$, então $\norm{x+y}^2 = \norm x^2 + \norm y^2$. Uma
família \omterm{def:b1:euclid:orthogonal}{ortogonal} de vetores não nulos é \omterm{def:b1:vspaces:free}{livre}. Numa \omterm{def:b1:vspaces:free}{base} \omterm{def:b1:euclid:orthogonal}{ortonormal}
$(e_1, \dots,
e_n)$, as \omterm{prop:b1:vspaces:coordinates}{coordenadas} e os produtos são
\[
x = \sum_{i} \langle x, e_i\rangle\, e_i,
\qquad
\langle x, y \rangle = \sum_i \langle x, e_i\rangle \langle y,
e_i\rangle,
\qquad
\norm x^2 = \sum_i \langle x, e_i\rangle^2 .
\]
\end{proposition}

\begin{proof}
Pitágoras: desenvolva. Liberdade: tome $\langle\,\cdot\,, x_j\rangle$ de uma combinação nula:
$\lambda_j \norm{x_j}^2 = 0$. \omterm{prop:b1:vspaces:coordinates}{Coordenadas}: escreva $x = \sum \lambda_i e_i$ e tome o produto com $e_j$: $\lambda_j = \langle x, e_j\rangle$;
as duas fórmulas seguem por bilinearidade.
\end{proof}

\begin{example}[Coordenadas ortonormais, com uma verificação de Parseval]
\label{ex:b1:euclid:parsevalcheck}
Desenvolva $x = (1, 2, 3)$ na \omterm{def:b1:vspaces:free}{base} \omterm{def:b1:euclid:orthogonal}{ortonormal} do \cref{exo:b1:euclid:3},
\[
e_1 = \tfrac{1}{\sqrt2}(1,1,0), \quad
e_2 = \tfrac{1}{\sqrt6}(1,-1,2), \quad
e_3 = \tfrac{1}{\sqrt3}(-1,1,1).
\]
Nenhum sistema a resolver --- três \omterm{def:b1:euclid:def}{produtos internos}:
\[
\langle x, e_1\rangle = \frac{3}{\sqrt2},
\qquad
\langle x, e_2\rangle = \frac{1 - 2 + 6}{\sqrt6} =
\frac{5}{\sqrt6},
\qquad
\langle x, e_3\rangle = \frac{-1 + 2 + 3}{\sqrt3} =
\frac{4}{\sqrt3}.
\]
Certificação pela fórmula da norma da \omterm{def:b1:logic:statement}{proposição}:
\[
\frac{9}{2} + \frac{25}{6} + \frac{16}{3}
= \frac{27 + 25 + 32}{6} = 14 = \norm{x}^2 = 1 + 4 + 9 .
\]
Esta verificação por soma dos quadrados das \omterm{prop:b1:vspaces:coordinates}{coordenadas} (uma
identidade de Parseval finita) custa segundos e pega erros de sinal
e de normalização com quase certeza --- torne-a um hábito sempre que
calcular um desenvolvimento \omterm{def:b1:euclid:orthogonal}{ortonormal}; a sua versão em dimensão
infinita, para os coeficientes de Fourier do \cref{ex:b1:euclid:trigortho}, é um teorema do
volume do terceiro ano de graduação.
\end{example}

\begin{theorem}[Gram--Schmidt]\label{thm:b1:euclid:gramschmidt}
Todo \omterm{def:b1:euclid:def}{espaço euclidiano} tem \omterm{def:b1:vspaces:free}{bases} \omterm{def:b1:euclid:orthogonal}{ortonormais}. Explicitamente, a
partir de qualquer \omterm{def:b1:vspaces:free}{base} $(v_1, \dots, v_n)$, a receita\index{processo de
Gram--Schmidt}
\[
w_k = v_k - \sum_{i=1}^{k-1} \langle v_k, e_i\rangle\, e_i ,
\qquad
e_k = \frac{w_k}{\norm{w_k}}
\]
produz uma \omterm{def:b1:vspaces:free}{base} \omterm{def:b1:euclid:orthogonal}{ortonormal} $(e_1, \dots, e_n)$ com $\operatorname{Vect}(e_1, \dots, e_k) = \operatorname{Vect}(v_1,
\dots, v_k)$ para todo $k$.
\end{theorem}

\begin{proof}
Indução sobre $k$. Supondo $(e_1, \dots, e_{k-1})$ \omterm{def:b1:euclid:orthogonal}{ortonormal} e gerando $\operatorname{Vect}(v_1, \dots, v_{k-1})$: o vetor
$w_k$ é \omterm{def:b1:euclid:orthogonal}{ortogonal} a cada $e_j$ ($j < k$) por construção ($\langle w_k, e_j \rangle = \langle v_k, e_j\rangle - \langle v_k,
e_j\rangle$), e
$w_k \neq 0$ pois $v_k \notin
\operatorname{Vect}(v_1, \dots, v_{k-1})$. Normalizar preserva a ortogonalidade; o enunciado
sobre o \omterm{def:b1:vspaces:span}{espaço gerado} vale porque $e_k$ é combinação de $v_k$ e dos
$e_i$ anteriores, de modo invertível.
\end{proof}

\begin{example}[Gram--Schmidt sobre polinômios, por extenso]
\label{ex:b1:euclid:legendre}
Ortonormalize $(1, X, X^2)$ em $\R_2[X]$ com $\langle P,
Q\rangle = \int_0^1 PQ$. \emph{Passo 1}: $\norm{1}^2 = 1$,
logo $e_1 = 1$. \emph{Passo 2}: $w_2 = X - \langle X, 1\rangle\,1 = X -
\frac12$, e $\norm{w_2}^2 = \int_0^1\bigl(x - \frac12\bigr)^2
\dd x = \frac1{12}$: $e_2 = \sqrt{12}\,\bigl(X - \frac12\bigr)$.
\emph{Passo 3}: $\langle X^2, e_1\rangle = \frac13$ e
\[
\langle X^2, e_2\rangle
= \sqrt{12}\int_0^1 x^2\Bigl(x - \frac12\Bigr)\dd x
= \frac{\sqrt{12}}{12},
\qquad\text{logo}\qquad
w_3 = X^2 - \frac13 - \Bigl(X - \frac12\Bigr)
= X^2 - X + \frac16 .
\]
A sua norma foi calculada no \cref{exo:b1:euclid:9}: $\norm{w_3}^2 =
\frac1{180}$, donde $e_3 = \sqrt{180}\,\bigl(X^2 - X +
\frac16\bigr)$. Os
\omterm{def:b1:poly:def}{polinômios} $1$, $X - \frac12$, $X^2 - X +
\frac16$ são, a menos de escala, os primeiros
\emph{\omterm{def:b1:poly:def}{polinômios} de Legendre} do \omterm{prop:b1:reals:intervals}{intervalo} $\intcc{0}{1}$; a construção
prossegue um grau de cada vez, cada novo \omterm{def:b1:poly:def}{polinômio} \omterm{def:b1:euclid:orthogonal}{ortogonal} a todos
os seus antecessores. Note como o algoritmo reaproveita o trabalho
anterior: a \omterm{def:b1:linmaps:projection}{projeção} subtraída no passo 3 é exatamente a melhor
aproximação afim de $X^2$ encontrada no \cref{ex:b1:euclid:bestapprox} --- Gram--Schmidt
\emph{é} \omterm{thm:b1:euclid:projection}{projeção ortogonal} iterada.
\end{example}

\begin{theorem}[Projeção ortogonal]\label{thm:b1:euclid:projection}
Seja $F$ um \omterm{def:b1:vspaces:subspace}{subespaço} do \omterm{def:b1:euclid:def}{espaço euclidiano} $E$. Então
\[
E = F \oplus F^{\perp},
\]
e a \omterm{def:b1:linmaps:projection}{projeção} associada $p_F$ sobre $F$ (a \emph{projeção
ortogonal}\index{projeção ortogonal}) é dada, em qualquer \omterm{def:b1:vspaces:free}{base}
\omterm{def:b1:euclid:orthogonal}{ortonormal} $(e_1, \dots, e_k)$ de $F$, por $p_F(x) = \sum_i
\langle x, e_i\rangle e_i$. Ela realiza a distância a $F$:
para todo $y \in F$,
\[
\norm{x - p_F(x)} \leq \norm{x - y},
\]
com igualdade apenas para $y = p_F(x)$; escreve-se $d(x, F) = \norm{x -
p_F(x)}$.
\end{theorem}

\begin{proof}
Tome uma \omterm{def:b1:vspaces:free}{base} \omterm{def:b1:euclid:orthogonal}{ortonormal} $(e_i)_{i \leq k}$ de $F$ (o \cref{thm:b1:euclid:gramschmidt} dentro de $F$) e
ponha $\pi(x) =
\sum \langle x, e_i\rangle e_i \in F$. Então $x - \pi(x) \perp e_j$ para cada $j$ (o mesmo cancelamento de
antes), logo $x - \pi(x) \in
F^\perp$: $E = F + F^\perp$. E $F \cap F^\perp = \{0\}$: tal vetor satisfaz $\langle x, x\rangle = 0$. Logo a
soma é direta e $\pi = p_F$.

Distância: para $y \in F$, decomponha $x - y = (x - p_F(x)) + (p_F(x) -
y)$, pedaços \omterm{def:b1:euclid:orthogonal}{ortogonais}
($F^\perp$ e $F$); Pitágoras:
\[
\norm{x - y}^2 = \norm{x - p_F(x)}^2 + \norm{p_F(x) - y}^2
\geq \norm{x - p_F(x)}^2,
\]
igualdade se e somente se $y = p_F(x)$.
\end{proof}

\begin{example}[As projeções nunca alongam]
\label{ex:b1:euclid:bessel}
Aplicando Pitágoras à decomposição $x = p_F(x) + (x - p_F(x))$:
\[
\norm{p_F(x)}^2 = \norm x^2 - \norm{x - p_F(x)}^2 \leq
\norm x^2 ,
\]
com igualdade se e somente se $x \in F$. Numa \omterm{def:b1:vspaces:free}{base} \omterm{def:b1:euclid:orthogonal}{ortonormal} $(e_1,
\dots, e_k)$
de $F$ isto lê-se $\sum_{i \leq k}\langle x,
e_i\rangle^2 \leq \norm x^2$ (uma \emph{desigualdade de Bessel}): por
mais direções \omterm{def:b1:euclid:orthogonal}{ortonormais} que se meçam, as \omterm{prop:b1:vspaces:coordinates}{coordenadas} ao quadrado
nunca ultrapassam o comprimento ao quadrado --- compare com a
igualdade exata do \cref{ex:b1:euclid:parsevalcheck} quando a família é uma \omterm{def:b1:vspaces:free}{base} completa. Esta
desigualdade de uma linha é o que torna somáveis os coeficientes de
Fourier no volume do terceiro ano de graduação; aqui ela já explica
por que acrescentar mais funções de \omterm{def:b1:vspaces:free}{base} a um ajuste de mínimos
quadrados só pode diminuir o resíduo.
\end{example}

\begin{example}[Melhor aproximação quadrática]\label{ex:b1:euclid:bestapprox}
Em $C(\intcc{0}{1})$ com $\langle f, g\rangle = \int_0^1 fg$, o \omterm{def:b1:poly:def}{polinômio} de grau $\leq 1$ mais próximo de
$f(x) = x^2$ na distância associada (em média quadrática) é $p_F(f)$, com
$F = \R_1[X]$. Gram--Schmidt sobre $(1, X)$: $e_1 = 1$, $w_2 = X - \frac12$, $\norm{w_2}^2 = \int_0^1 (x - \frac12)^2 = \frac{1}{12}$, $e_2 =
\sqrt{12}\,(X - \tfrac12)$.
Então
\[
p_F(f) = \langle f, e_1\rangle e_1 + \langle f, e_2\rangle e_2
= \frac13 + 12\Bigl(\int_0^1 x^2\bigl(x - \tfrac12\bigr)\dd
x\Bigr)\bigl(X - \tfrac12\bigr)
= X - \frac{1}{6},
\]
usando $\int_0^1 x^2(x - \frac12)\dd x = \frac14 - \frac16 =
\frac{1}{12}$. A ideia dos ``mínimos quadrados'' numa linha de
álgebra linear.
\end{example}

\begin{method}[Três vias para uma distância $d(x, F)$]
\label{met:b1:euclid:distance}
\begin{enumerate}
  \item \emph{\omterm{def:b1:vspaces:free}{Base} \omterm{def:b1:euclid:orthogonal}{ortonormal} de $F$}: então $p_F(x) = \sum_i
        \langle x, e_i\rangle e_i$ e, por Pitágoras,
        \[
        d(x, F)^2 = \norm{x}^2 - \norm{p_F(x)}^2
        = \norm x^2 - \sum_i \langle x, e_i\rangle^2 ,
        \]
        muitas vezes mais barato do que calcular o próprio $x - p_F(x)$.
  \item \emph{Equações normais}: com qualquer família geradora de
        $F$, resolva $\langle x - p, v_j\rangle = 0$ para os coeficientes de $p$ (o \cref{exo:b1:euclid:5}) ---
        sem nenhuma ortonormalização.
  \item \emph{Pelo complemento}: se $F^\perp$ é menor do que $F$
        (por exemplo, $F$ um \omterm{def:b1:linmaps:forms}{hiperplano} e $F^\perp$ uma reta
        $\operatorname{Vect}(n)$), projete sobre $F^\perp$ em vez disso:
        \[
        d(x, F) = \norm{p_{F^\perp}(x)} =
        \frac{\abs{\langle x, n\rangle}}{\norm n} ,
        \]
        que é a fórmula clássica da distância a um plano (o \cref{exo:b1:multivar:8} a
        usa).
\end{enumerate}
A via 3 é um caso particular de um reflexo geral: projete sempre
sobre aquele dentre $F$ e $F^\perp$ que tiver a menor dimensão.
\end{method}

\begin{example}[Ortogonalidade trigonométrica: uma prévia de Fourier]
\label{ex:b1:euclid:trigortho}
Em $C(\intcc{0}{2\pi})$ com $\langle f, g\rangle =
\frac1\pi\int_0^{2\pi} fg$, a família
\[
\Bigl(\frac{1}{\sqrt2},\ \cos x,\ \sin x,\ \cos 2x,\ \sin 2x,\
\dots\Bigr)
\]
é \omterm{def:b1:euclid:orthogonal}{ortonormal}: por exemplo $\langle \cos px, \cos qx\rangle =
\frac1\pi\int_0^{2\pi}\cos px\cos qx\,\dd x = 0$ para $p \neq q$ (linearize o produto em
$\frac12[\cos(p{-}q)x +
\cos(p{+}q)x]$ e integre sobre períodos completos), enquanto $\frac1\pi\int_0^{2\pi}\cos^2 px\,\dd x = 1$. A
\omterm{thm:b1:euclid:projection}{projeção ortogonal} sobre o \omterm{def:b1:vspaces:span}{espaço gerado} pelas primeiras $2N + 1$
destas funções tem portanto \omterm{prop:b1:vspaces:coordinates}{coordenadas} $\langle f, e_i\rangle$ --- integrais contra
cossenos e senos. Estes são os \emph{coeficientes de Fourier} de
$f$, e a \omterm{def:b1:linmaps:projection}{projeção} é a sua melhor aproximação trigonométrica em média
quadrática; o volume do terceiro ano de graduação estuda a sua
convergência. A ortogonalidade faz todo o trabalho: as fórmulas dos
coeficientes são o \cref{thm:b1:euclid:projection}, palavra por palavra.
\end{example}

\section{Isometrias do plano}

\begin{definition}\label{def:b1:euclid:isometry}
Um endomorfismo $u$ de um \omterm{def:b1:euclid:def}{espaço euclidiano} é uma
\emph{isometria}\index{isometria} (ou \omterm{def:b1:logic:map}{aplicação} \omterm{def:b1:euclid:orthogonal}{ortogonal}) quando
preserva a norma: $\norm{u(x)} = \norm x$ para todo $x$ --- equivalentemente (por
polarização), quando preserva o \omterm{def:b1:euclid:def}{produto interno}; equivalentemente,
quando a sua matriz $A$ numa \omterm{def:b1:vspaces:free}{base} \omterm{def:b1:euclid:orthogonal}{ortonormal} satisfaz $A^{\mathsf T} A = I$. As
isometrias formam um \omterm{def:b1:structures:group}{grupo}, o \omterm{def:b1:structures:group}{grupo} \omterm{def:b1:euclid:orthogonal}{ortogonal} $O(E)$.
\end{definition}

\begin{example}[Reconhecer uma isometria de relance]
\label{ex:b1:euclid:recognize}
Será $A = \dfrac15\begin{pmatrix} 3 & -4\\ 4 & 3\end{pmatrix}$ \omterm{def:b1:euclid:orthogonal}{ortogonal}? Colunas: normas $\frac15\sqrt{9 + 16} = 1$ e $\frac15\sqrt{16 + 9} = 1$; produto $\frac1{25}(3\cdot(-4) +
4\cdot3) = 0$.
Sim --- e $\det A = \frac{9 + 16}{25} = 1$, logo ela é a rotação $R_\theta$ com $\cos\theta = \frac35$, $\sin\theta = \frac45$ (a
``rotação $3$-$4$-$5$'', cujo ângulo não é fração notável alguma de
$\pi$). Em contraste, $B =
\frac{1}{\sqrt2}\begin{pmatrix} 1 & 1\\ 0 & 1\end{pmatrix}$ tem aparência de determinante unitário
mas primeira coluna de norma diferente de $1$ ($\frac1{\sqrt2}$): não é
\omterm{def:b1:euclid:orthogonal}{ortogonal} --- determinante $\pm1$ sozinho não certifica nada, as
colunas têm de ser verificadas.
\end{example}

\begin{theorem}[Isometrias do plano]\label{thm:b1:euclid:plane}
Numa \omterm{def:b1:vspaces:free}{base} \omterm{def:b1:euclid:orthogonal}{ortonormal} de um plano euclidiano, as matrizes das
\omterm{def:b1:euclid:isometry}{isometrias} são exatamente
\[
R_\theta = \begin{pmatrix}
\cos\theta & -\sin\theta\\ \sin\theta & \cos\theta
\end{pmatrix}
\quad (\text{rotação de ângulo } \theta,\ \det = 1),
\]
\[
S_\theta = \begin{pmatrix}
\cos\theta & \sin\theta\\ \sin\theta & -\cos\theta
\end{pmatrix}
\quad (\det = -1),
\]
sendo a última a reflexão na reta que faz ângulo $\frac\theta2$ com o
primeiro vetor da \omterm{def:b1:vspaces:free}{base}.
\end{theorem}

\begin{proof}
Seja $A = \begin{pmatrix} a & c\\ b & d\end{pmatrix}$ com $A^{\mathsf T}A = I$: as colunas são unitárias e \omterm{def:b1:euclid:orthogonal}{ortogonais}. A
primeira coluna é $(\cos\theta, \sin\theta)$ para algum $\theta$; a segunda, unitária e
\omterm{def:b1:euclid:orthogonal}{ortogonal} a ela, é $\pm(-\sin\theta, \cos\theta)$. O sinal $+$ dá $R_\theta$; o sinal $-$ dá
$S_\theta$. Verifica-se que $S_\theta^2 = I$ e que o vetor $(\cos\frac\theta2,
\sin\frac\theta2)$ é fixo enquanto o
seu \omterm{def:b1:euclid:orthogonal}{ortogonal} é invertido: uma reflexão. (E $R_\alpha R_\beta = R_{\alpha+\beta}$: o \omterm{def:b1:structures:group}{grupo} das
rotações é o \omterm{def:b1:structures:group}{grupo} dos ângulos --- compare com o \cref{thm:b1:complex:funceq}.)
\end{proof}

\begin{omfigure}
\begin{tikzpicture}[scale=1.15]
  \draw[->, thin] (-1.2,0) -- (2.9,0) node[below] {$x$};
  \draw[->, thin] (0,-1.2) -- (0,2.6) node[left] {$y$};
  \draw[omDef, very thick] (-0.9,0) -- (2.6,0)
    node[below, pos=0.9, font=\small] {eixo de $r_1$};
  \draw[omThm, very thick] (-0.9,-0.9) -- (2.3,2.3)
    node[right, pos=0.97, font=\small] {eixo de $r_2$};
  \fill (2,0.5) circle (0.05) node[right, font=\small] {$M$};
  \fill (2,-0.5) circle (0.05)
    node[right, font=\small] {$r_1(M)$};
  \fill (-0.5,2) circle (0.05)
    node[above, font=\small] {$r_2(r_1(M))$};
  \draw[omDef, dashed] (2,0.5) -- (2,-0.5);
  \draw[omThm, dashed] (2,-0.5) -- (-0.5,2);
  \draw[omProp, thick, ->] (0.9,0) arc (0:45:0.9);
  \node[omProp, font=\small] at (1.15,0.32) {$\frac\pi4$};
\end{tikzpicture}

{\small Duas reflexões fazem uma rotação: refletir $M = (2,
0.5)$ no eixo $x$
e depois na reta $y = x$ leva a $(-0.5, 2)$ --- a \omterm{def:b1:linmaps:kerim}{imagem} de $M$ pela
rotação de ângulo $\frac\pi2$ em torno da origem, o \emph{dobro} do ângulo
$\frac\pi4$ entre os eixos. O problema de fim de semana transforma esta
figura na \omterm{def:b1:structures:law}{lei de composição} de todas as \omterm{def:b1:euclid:isometry}{isometrias} do plano.}
\end{omfigure}

\begin{remark}[Armadilhas comuns]
\label{rem:b1:euclid:pitfalls}
\emph{A fórmula da \omterm{def:b1:linmaps:projection}{projeção} precisa de uma \omterm{def:b1:vspaces:free}{base} \omterm{def:b1:euclid:orthogonal}{ortonormal}}: para
uma família $(v_i)$ apenas geradora de $F$, a soma $\sum_i\langle
x, v_i\rangle v_i$ \emph{não}
é $p_F(x)$ (teste $F = \R^2$, $v_1
= e_1$, $v_2 = e_1 + e_2$); com uma família não
\omterm{def:b1:euclid:orthogonal}{ortonormal}, resolva antes as equações normais (o \cref{met:b1:euclid:distance}~(2)).
\emph{As famílias \omterm{def:b1:euclid:orthogonal}{ortogonais} têm de evitar o $0$ para serem \omterm{def:b1:vspaces:free}{livres}}:
o vetor nulo é \omterm{def:b1:euclid:orthogonal}{ortogonal} a tudo, inclusive a si mesmo --- a
liberdade na \cref{prop:b1:euclid:orthfree} exige vetores não nulos.
\emph{$F^\perp$ depende do \omterm{def:b1:euclid:def}{produto interno}}: em $\R_1[X]$, o
complemento de $\operatorname{Vect}(X)$ para $\int_0^1 PQ$ não são as constantes, mas $\operatorname{Vect}(1 - \frac32 X)$
--- calcule $\int_0^1 x(1 - \frac32 x) = \frac12 - \frac12 = 0$; ``perpendicular'' não faz sentido enquanto o
produto não for nomeado.
\emph{Não desenvolva $\norm{x + y}$ linearmente}: a identidade correta é
$\norm{x+y}^2 = \norm x^2 + 2\langle x, y\rangle +
\norm y^2$; o termo cruzado só se anula sob ortogonalidade (Pitágoras),
e a desigualdade triangular é uma \emph{desigualdade}.
\emph{Levar vetores unitários em vetores unitários não basta}:
$u(x,
y) = (x + y,\ 0)$ leva os dois vetores da \omterm{def:b1:vspaces:free}{base} canônica no vetor unitário
$(1, 0)$, e no entanto $\norm{u(1,1)} = 2 \neq \sqrt2$: nenhuma \omterm{def:b1:euclid:isometry}{isometria}. A definição exige
$\norm{u(x)} = \norm x$ para \emph{todo} $x$; em termos matriciais $A^{\mathsf T}A = I$, isto é,
colunas unitárias \emph{e dois a dois \omterm{def:b1:euclid:orthogonal}{ortogonais}} --- as duas
condições, verificadas juntas.
\end{remark}

\begin{remark}[Para onde vai o produto interno]
\label{rem:b1:euclid:whereused}
A \omterm{thm:b1:euclid:projection}{projeção ortogonal} é o teorema mais aplicado do capítulo: ela
sustenta os mínimos quadrados (o problema de fim de semana do \cref{ch:b1:multivar}
constrói retas de regressão sobre ela), os coeficientes de Fourier
(o \cref{ex:b1:euclid:trigortho}) e as equações normais do \cref{exo:b1:euclid:5}, que a análise numérica
resolve em grande escala. A classificação $R_\theta / S_\theta$ é completada abaixo:
o problema de fim de semana classifica \emph{todas} as
transformações do plano que preservam distâncias, \omterm{def:b1:linmaps:def}{lineares} ou não, e
os seus \omterm{def:b1:structures:group}{grupos} finitos --- a matemática das rosetas e dos polígonos
regulares. No volume do segundo ano de graduação o \omterm{def:b1:euclid:def}{produto interno}
encontra a teoria dos autovalores (matrizes simétricas, formas
quadráticas); no terceiro ano, a geometria euclidiana em dimensão
infinita torna-se a teoria dos espaços de Hilbert.
\end{remark}

\begin{remark}[Perspectivas dentro do Livro 3]
\label{rem:b1:euclid:perspectives}
Duas pontes saem deste capítulo. \emph{Para trás}, rumo à álgebra
linear: a matriz de Gram do \cref{exo:b1:euclid:11} empacota \omterm{def:b1:euclid:def}{produtos internos} na
maquinaria de determinantes do \cref{ch:b1:det}, e a \omterm{thm:b1:euclid:projection}{projeção ortogonal} é o
projetor especial do \cref{ch:b1:linmaps} cujo núcleo é $F^\perp$ --- toda a sua
álgebra ($p^2 = p$, $s = 2p -
\mathrm{id}$) aplica-se palavra por palavra, agora com
o bônus de que $\norm{x
- p(x)}$ é uma \emph{distância}. \emph{Para a frente},
rumo à análise: o \cref{ch:b1:curves} mede comprimentos de arco com a norma deste
capítulo e não classifica nada sem as suas \omterm{def:b1:euclid:isometry}{isometrias}; o \cref{ch:b1:multivar} lê o
gradiente através de Cauchy--Schwarz (a subida mais íngreme) e fecha
o volume com os mínimos quadrados, que são o \cref{thm:b1:euclid:projection} aplicado a um
vetor de dados em $\R^n$. O \omterm{def:b1:euclid:def}{produto interno} é o ponto em que a
álgebra e a análise do livro finalmente se encontram.
\end{remark}

\section{Exercícios}

\begin{exercise}[$\star$]\label{exo:b1:euclid:1}
Em $\R^3$ canônico: calcule $\langle u, v\rangle$, $\norm u$, $\norm v$ e o ângulo
entre $u = (1, 2, 2)$ e $v = (2, -2,
1)$. Verifique Cauchy--Schwarz numericamente.
\end{exercise}

\begin{exercise}[$\star$]\label{exo:b1:euclid:2}
Demonstre a identidade do paralelogramo $\norm{x+y}^2 + \norm{x-y}^2 =
2\norm x^2 + 2\norm y^2$ em qualquer \omterm{def:b1:euclid:def}{espaço
euclidiano}, e use-a para mostrar que a norma do \omterm{def:b1:reals:bounds}{supremo} em $\R^2$,
$\norm{(x,y)}_\infty = \max(\abs x,
\abs y)$, não provém de um \omterm{def:b1:euclid:def}{produto interno}.
\end{exercise}

\begin{exercise}[$\star$]\label{exo:b1:euclid:3}
Aplique Gram--Schmidt a $\bigl((1,1,0), (1,0,1), (0,1,1)\bigr)$ no $\R^3$ canônico.
\end{exercise}

\begin{exercise}[$\star$]\label{exo:b1:euclid:4}
Em $\R^3$, seja $F = \operatorname{Vect}\bigl((1,1,1)\bigr)$. Determine $F^\perp$ (equação e \omterm{def:b1:vspaces:free}{base}), a
matriz de $p_F$ na \omterm{def:b1:vspaces:free}{base} canônica, e $d\bigl((1, 2, 3), F\bigr)$.
\end{exercise}

\begin{exercise}[$\star\star$]\label{exo:b1:euclid:5}
(Equações normais) Sejam $F =
\operatorname{Vect}\bigl((1,0,1),(0,1,1)\bigr) \subseteq \R^3$ e $x = (1, 1, 4)$. Calcule $p_F(x)$ resolvendo
$\langle x - p, v
\rangle = 0$ para os dois geradores ($p = \alpha(1,0,1) +
\beta(0,1,1)$), e depois $d(x, F)$. Por que
Gram--Schmidt é desnecessário aqui?
\end{exercise}

\begin{exercise}[$\star\star$]\label{exo:b1:euclid:6}
Para $f$ \omterm{def:b1:continuity:continuous}{contínua} em $\intcc{0}{1}$, demonstre
\[
\Bigl(\int_0^1 f\Bigr)^{\!2} \leq \int_0^1 f^2 ,
\]
com o caso de igualdade, como instância de Cauchy--Schwarz em $C(\intcc{0}{1})$.
Demonstre em seguida $\bigl(\sum_{i=1}^n a_i\bigr)^2 \leq
n \sum a_i^2$ para reais $a_i$.
\end{exercise}

\begin{exercise}[$\star\star$]\label{exo:b1:euclid:7}
Demonstre que, para todo \omterm{def:b1:vspaces:subspace}{subespaço} $F$ de um \omterm{def:b1:euclid:def}{espaço euclidiano}:
$(F^{\perp})^{\perp} = F$ e $\dim F^{\perp} = \dim E - \dim F$.
\end{exercise}

\begin{exercise}[$\star\star$]\label{exo:b1:euclid:8}
Identifique as \omterm{def:b1:euclid:isometry}{isometrias} do plano de matrizes
\[
A = \frac{1}{\sqrt 2}\begin{pmatrix} 1 & -1\\ 1 & 1\end{pmatrix},
\qquad
B = \frac{1}{5}\begin{pmatrix} 3 & 4\\ 4 & -3\end{pmatrix}
\]
(tipo, ângulo ou eixo). Calcule $A^8$ e $B^2$ sem multiplicar
matrizes.
\end{exercise}

\begin{exercise}[$\star\star\star$]\label{exo:b1:euclid:9}
(Mínimo como \omterm{def:b1:linmaps:projection}{projeção}) Calcule
\[
\min_{(a, b) \in \R^2} \int_0^1 \bigl(x^2 - a - bx\bigr)^2 \dd x ,
\]
usando o \cref{ex:b1:euclid:bestapprox}: o mínimo é $\norm{f -
p_F(f)}^2$ para $f = X^2$, $F = \R_1[X]$.
\end{exercise}

\begin{exercise}[$\star\star\star$]\label{exo:b1:euclid:10}
Seja $u$ uma \omterm{def:b1:euclid:isometry}{isometria} de um \omterm{def:b1:euclid:def}{espaço euclidiano} $E$. Demonstre que
$\ker(u - \mathrm{id}) \perp \operatorname{im}(u - \mathrm{id})$, e deduza $E = \ker(u - \mathrm{id}) \oplus \operatorname{im}(u -
\mathrm{id})$. \emph{(Calcule $\langle x - u(x), y\rangle$ para $u(y) = y$, usando a
preservação do produto.)}
\end{exercise}

\begin{exercise}[$\star\star$]\label{exo:b1:euclid:11}
(Matriz de Gram) Para vetores $v_1, \dots, v_k$ de um \omterm{def:b1:euclid:def}{espaço euclidiano}, seja
$G = \bigl(\langle v_i, v_j\rangle\bigr)_{1 \leq i, j
\leq k}$ a sua \emph{matriz de Gram}.
\begin{enumerate}
  \item Demonstre que $(v_1, \dots, v_k)$ é \omterm{def:b1:vspaces:free}{livre} se e somente se $G$ é
        invertível. \emph{(Se $Gc = 0$, calcule $\norm{\sum_i c_i v_i}^2$.)}
  \item Calcule a matriz de Gram de $(1, X, X^2)$ em $\R_2[X]$ com $\langle P, Q\rangle = \int_0^1 PQ$,
        reconheça a \omterm{pb:b1:det:1}{matriz de Hilbert} $H_3$ do problema de fim de
        semana do \cref{ch:b1:det}, e conclua a liberdade a partir de $\det H_3 =
        \frac1{2160} \neq 0$.
\end{enumerate}
\end{exercise}

\begin{exercise}[$\star\star\star$]\label{exo:b1:euclid:12}
Seja $u$ um endomorfismo de um \omterm{def:b1:euclid:def}{espaço euclidiano} $E$ cuja matriz
$A$ numa \omterm{def:b1:vspaces:free}{base} \omterm{def:b1:euclid:orthogonal}{ortonormal} é ao mesmo tempo \omterm{def:b1:euclid:orthogonal}{ortogonal} ($A^{\mathsf T}A
= I$) e
simétrica ($A^{\mathsf T} = A$).
\begin{enumerate}
  \item Mostre que $A^2 = I$, e deduza (via o \cref{thm:b1:linmaps:projchar}) que $E = \ker(u -
        \mathrm{id}) \oplus \ker(u + \mathrm{id})$.
  \item Mostre que os dois \omterm{def:b1:vspaces:subspace}{subespaços} são \omterm{def:b1:euclid:orthogonal}{ortogonais}, de modo que
        $u$ é a \emph{simetria \omterm{def:b1:euclid:orthogonal}{ortogonal}} em relação a $F =
        \ker(u - \mathrm{id})$: a
        reflexão através de $F$. \emph{(Para $u(x) = x$ e $u(y) = -y$, calcule
        $\langle
        x, y\rangle$ de duas maneiras.)}
  \item Classifique o caso plano: quais matrizes do \cref{thm:b1:euclid:plane} são
        simétricas, e quais são as aplicações correspondentes?
\end{enumerate}
\end{exercise}

\section{Problema: isometrias do plano e o teorema de Leonardo}

\begin{problem}\label{pb:b1:euclid:1}
Uma \emph{\omterm{def:b1:euclid:isometry}{isometria} do plano} é qualquer \omterm{def:b1:logic:map}{aplicação} $f \colon \R^2 \to
\R^2$ que
preserva distâncias: $\norm{f(x) - f(y)} = \norm{x - y}$ para todos $x, y$ --- sem nenhuma
hipótese de linearidade. Este problema demonstra que tais aplicações
são exatamente as translações, as rotações, as reflexões e as
\omterm{pb:b1:euclid:1}{reflexões deslizantes} (a classificação das \omterm{def:b1:euclid:isometry}{isometrias} do plano),
calcula as suas composições e determina todos os seus \omterm{def:b1:structures:group}{grupos}
\emph{finitos}: o teorema de Leonardo, a matemática por trás dos
padrões em roseta. A partir da Parte II identificamos $\R^2$ com
$\C$ (o \cref{ch:b1:complex}): o \omterm{def:b1:euclid:def}{produto interno} canônico é $\langle z, w\rangle = \operatorname{Re}(z\conj w)$ e a norma é o
\omterm{def:b1:complex:field}{módulo}.\index{isometria}\index{reflexão deslizante}
\index{teorema de Leonardo}\index{grupo diedral}

\medskip
\noindent\textbf{Parte I --- Toda \omterm{def:b1:euclid:isometry}{isometria} é afim.}
\begin{enumerate}
  \item Verifique que as translações $t_a(x) = x + a$, as \omterm{def:b1:euclid:isometry}{isometrias} \omterm{def:b1:linmaps:def}{lineares}
        e todas as suas composições são \omterm{def:b1:euclid:isometry}{isometrias}, e que as
        \omterm{def:b1:euclid:isometry}{isometrias} formam um \omterm{def:b1:structures:group}{grupo} sob composição.
  \item Seja $f$ uma \omterm{def:b1:euclid:isometry}{isometria} com $f(0) = 0$. Mostre que $f$ preserva
        normas e, em seguida --- por polarização, o \cref{thm:b1:euclid:cs} ---, que
        $\langle f(x),
        f(y)\rangle = \langle x, y\rangle$ para todos $x, y$.
  \item Ainda com $f(0) = 0$: desenvolva $\norm{f(x + y) - f(x) -
        f(y)}^2$ e $\norm{f(\lambda x) - \lambda f(x)}^2$ usando a questão
        2, e conclua que $f$ é \emph{\omterm{def:b1:linmaps:def}{linear}}: $f \in O(\R^2)$.
  \item Deduza que toda \omterm{def:b1:euclid:isometry}{isometria} $f$ se escreve de modo
        \emph{único} como $f = t_a \circ g$ com $a = f(0)$ e $g$ uma \omterm{def:b1:euclid:isometry}{isometria}
        \omterm{def:b1:linmaps:def}{linear} (a \emph{parte \omterm{def:b1:linmaps:def}{linear}} de $f$).
  \item Chame $f$ de \emph{direta} se $\det g = 1$, e de \emph{indireta}
        se $\det g = -1$. Mostre que a parte \omterm{def:b1:linmaps:def}{linear} de uma composição é a
        composição das partes \omterm{def:b1:linmaps:def}{lineares}, e enuncie a regra de sinais
        resultante (direta/indireta compõem-se como $+1/-1$).
\end{enumerate}

\medskip
\noindent\textbf{Parte II --- Os quatro tipos.} Pelo \cref{thm:b1:euclid:plane}, as
\omterm{def:b1:euclid:isometry}{isometrias} \omterm{def:b1:linmaps:def}{lineares} de $\C$ são $z
\mapsto az$ e $z \mapsto a\conj z$ com $\abs a = 1$; logo toda
\omterm{def:b1:euclid:isometry}{isometria} do plano é
\[
f(z) = a z + b \quad (\text{direta}) \qquad\text{ou}\qquad
f(z) = a\conj z + b \quad (\text{indireta}), \qquad \abs a = 1 .
\]
\begin{enumerate}[resume]
  \item Verifique o dicionário: $R_\theta$ é $z \mapsto
        \eu^{\iu\theta}z$ e $S_\theta$ é $z \mapsto
        \eu^{\iu\theta}\conj z$
        (confira ambos em $1$ e em $\iu$).
  \item (Caso direto) Seja $f(z) = az + b$, $\abs a = 1$. Mostre: se $a = 1$, $f$ é
        uma translação; se $a \neq 1$, $f$ tem o único ponto fixo $z_0 = b/(1 - a)$ e
        $f(z) - z_0
        = a(z - z_0)$: uma \emph{rotação} de centro $z_0$ e ângulo $\arg a$.
  \item (Caso indireto) Seja $f(z) = a\conj z + b$, e $v =
        a\conj b + b$. Mostre que $f \circ f = t_v$ e
        $f \circ t_v =
        t_v \circ f$. Se $v = 0$: mostre que o ponto médio de $z$ e $f(z)$
        é um ponto fixo, e que $f$ é uma \emph{reflexão} numa reta.
        Se $v \neq 0$: mostre que $r
        = t_{-v/2}\circ f$ é uma reflexão cujo eixo é
        paralelo a $v$, de modo que $f = t_{v/2} \circ r$ é uma \emph{\omterm{pb:b1:euclid:1}{reflexão
        deslizante}}. Conclua: toda \omterm{def:b1:euclid:isometry}{isometria} do plano é uma
        translação, uma rotação, uma reflexão ou uma \omterm{pb:b1:euclid:1}{reflexão
        deslizante} (a classificação das \omterm{def:b1:euclid:isometry}{isometrias} do plano).
  \item (Composições) Mostre: a composição de rotações de ângulos
        $\alpha$ e $\beta$ é uma rotação de ângulo $\alpha + \beta$ (uma
        translação se $\alpha + \beta \in
        2\pi\Z$); a composição de duas reflexões é uma
        rotação de ângulo igual ao dobro do ângulo entre os eixos
        (uma translação se os eixos são paralelos).
  \item Deduza que toda \omterm{def:b1:euclid:isometry}{isometria} do plano é composição de no
        máximo três reflexões.
\end{enumerate}

\medskip
\noindent\textbf{Parte III --- Três identificações.}
\begin{enumerate}[resume]
  \item Classifique completamente $f(z) = \iu\conj z + 1 + \iu$: tipo, eixo, vetor de
        deslizamento.
  \item Seja $f$ a rotação de ângulo $\frac\pi2$ em torno de $0$ e $g$ a
        rotação de ângulo $\frac\pi2$ em torno de $1$. Calcule $g \circ f$ na
        forma $z \mapsto az + b$ e identifique-a (tipo, centro, ângulo).
  \item Sejam $r_1(z) = \conj z$ (reflexão no eixo real) e $r_2(z) = \iu\conj z$ (reflexão na
        reta $y = x$). Calcule $r_2 \circ r_1$ e confira a questão 9 neste
        exemplo.
\end{enumerate}

\medskip
\noindent\textbf{Parte IV --- \omterm{def:b1:structures:group}{Grupos} finitos: o teorema de
Leonardo.} Seja $G$ um \omterm{def:b1:structures:group}{grupo} \emph{finito} de \omterm{def:b1:euclid:isometry}{isometrias} do plano.
\begin{enumerate}[resume]
  \item Mostre que as \omterm{def:b1:euclid:isometry}{isometrias} preservam baricentros: se $\lambda_1 + \dots + \lambda_m = 1$ e
        $f = t_a \circ
        g$ ($g$ \omterm{def:b1:linmaps:def}{linear}), então $f\bigl(\sum_i \lambda_i
        x_i\bigr) = \sum_i \lambda_i f(x_i)$.
  \item Ponha $c = \frac{1}{\abs G}\sum_{g \in G} g(x_0)$ para um $x_0$ escolhido qualquer. Mostre que
        todo $h \in G$ fixa $c$: um \omterm{def:b1:structures:group}{grupo} finito de \omterm{def:b1:euclid:isometry}{isometrias} tem
        um ponto fixo comum.
  \item Deduza que, depois de conjugar por $t_{-c}$, pode-se supor
        $G \subseteq O(\R^2)$. Seja $G^{+} = \{g \in G :
        \det g = 1\}$; mostre que ou $G = G^{+}$, ou $G^{+}$ tem
        exatamente índice $2$ em $G$ (exiba uma bijeção $G^+
        \to G \setminus G^+$).
  \item Mostre que um \omterm{def:b1:structures:group}{grupo} finito de \emph{rotações} em torno de
        $c$ é cíclico: entre os seus elementos escolha a rotação
        $R_{\theta_0}$ de menor ângulo $\theta_0 \in
        \intoo{0}{2\pi}$, e use a divisão euclidiana de
        ângulos para demonstrar que ela gera o \omterm{def:b1:structures:group}{grupo}; conclua $\theta_0 =
        \frac{2\pi}n$
        e $G^{+} = \{R_{\theta_0}^{\,k}\} \cong
        C_n$.
  \item Suponha $G \neq G^{+}$ e escolha uma reflexão $s \in G$. Mostre que
        $G = G^{+} \cup sG^{+}$, que $s r s = r^{-1}$ para toda rotação $r \in G^{+}$, e que todos os
        $n$ elementos de $sG^{+}$ são reflexões: $G$ é o \emph{\omterm{pb:b1:euclid:1}{grupo
        diedral}} $D_n$, de ordem $2n$.
  \item Conclua (\emph{teorema de Leonardo}): todo \omterm{def:b1:structures:group}{grupo} finito de
        \omterm{def:b1:euclid:isometry}{isometrias} do plano é cíclico $C_n$ ou diedral $D_n$.
\end{enumerate}

\medskip
\noindent\textbf{Parte V --- Dividendos, e síntese.}
\begin{enumerate}[resume]
  \item Mostre diretamente que um \omterm{def:b1:structures:group}{grupo} finito de \omterm{def:b1:euclid:isometry}{isometrias} não
        pode conter nenhuma translação nem nenhuma \omterm{pb:b1:euclid:1}{reflexão
        deslizante} além da identidade (considere as potências de
        tal elemento).
  \item Seja $P_n$ o $n$-ágono regular cujos vértices são as raízes
        $n$-ésimas da unidade ($n \geq 3$). Mostre que o seu \omterm{def:b1:structures:group}{grupo} de
        simetrias é exatamente $D_n$: as $n$ rotações $z
        \mapsto \omega^k z$ e as $n$
        reflexões $z \mapsto
        \omega^k \conj z$, $\omega = \eu^{2\iu\pi/n}$, e nenhuma outra.
  \item Exiba figuras planas cujos \omterm{def:b1:structures:group}{grupos} de simetria sejam,
        respectivamente, $C_1$, $D_1$, $D_2$ e $C_3$.
  \item Liste os oito elementos do \omterm{def:b1:structures:group}{grupo} de simetrias do quadrado de
        vértices $\pm1, \pm\iu$ como aplicações $z \mapsto
        \omega^k z$ ou $z \mapsto \omega^k\conj z$, e dê o eixo
        de cada uma das quatro reflexões.
  \item Sejam $f, g$ rotações do mesmo ângulo $\theta
        \notin 2\pi\Z$ em torno de
        centros \emph{distintos} $c_1 \neq
        c_2$. Calcule $f\circ g - g\circ f$ ponto a ponto
        e mostre que $f\circ g \neq g\circ f$; mostre além disso que $(f \circ
        g)\circ(g\circ f)^{-1}$ é uma
        translação não trivial, de modo que todo \omterm{def:b1:structures:group}{grupo} que contenha
        $f$ e $g$ é infinito --- uma segunda explicação para o
        centro único no teorema de Leonardo.
  \item Síntese, em quatro frases: quais dois resultados
        estruturais reduzem \omterm{def:b1:euclid:isometry}{isometrias} arbitrárias à álgebra linear
        (questões 3--4) e \omterm{def:b1:structures:group}{grupos} finitos arbitrários a \omterm{def:b1:structures:subgroup}{subgrupos} de
        $O(2)$ (questão 15); qual é a lista completa das \omterm{def:b1:euclid:isometry}{isometrias}
        do plano e quais invariantes (direta/indireta, pontos fixos)
        separam os quatro tipos; por que as regras de composição da
        questão 9 fazem das reflexões os geradores de tudo; e o que
        o teorema de Leonardo acrescenta na escala finita. Nomeie os
        dois teoremas demonstrados nas Partes II e IV.
\end{enumerate}
\end{problem}
